\mwnewpart
\renewcommand\runningtitle{Section 7. Problem Solving}
\chapter{Problem Solving}
\sectioncontents

\section{If the Computer Does Not Start}\ixs{startup problems}

\begin{longopts}
\lopt{\parbox[t]{\linewidth}{\ttfamily Non-System disk or disk error\\Replace and strike any key when ready}}{The
  diskette in drive A does not hold the RetrOS/LT system files. Put in the
  DOS diskette, or one made with \typed{FORMAT /S}, and press a key.}
\lopt{Disk Boot failure}{The startup disk could not be read. Try another
  copy of the diskette; clean the drive heads.}
\lopt{\parbox[t]{\linewidth}{\ttfamily Missing RETDOS.COM\\Error loading RETDOS.COM}}{\file{RETDOS.COM} is not
  the second file of the root directory, or could not be read. Put the
  system on the disk again with \typed{SYS}, from the DOS diskette.}
\lopt{Bad or missing Command Interpreter}{\file{COMMAND.COM} is not in the
  root directory of the startup disk, or is from another version of DOS.
  Copy it there from the DOS diskette.}
\lopt{\parbox[t]{\linewidth}{\ttfamily Invalid partition table\\Error loading operating system\\Missing operating system}}{The
  computer was starting from the fixed disk. The partition table has more
  than one active partition, the DOS partition could not be read, or it
  holds no system. Start from the DOS diskette, and use \typed{FDISK} to
  look at the partitions, or \typed{SYS C:} to put the system back.}
\lopt{Bad or missing \var{file}}{A \typed{DEVICE=} line in \file{CONFIG.SYS}
  names a file that is not there. Correct the line.}
\end{longopts}

If the computer starts ROM BASIC instead, it found no startup diskette in
drive A and no active partition on the fixed disk.

\section{Device Errors}\ixs{Abort, Retry, Ignore}

When a disk or a printer cannot do what it was asked, RetrOS/LT says what
happened and asks what to do:

\begin{dialogue}
Not ready error reading drive A
Abort, Retry, Ignore?
\end{dialogue}

\begin{tabularx}{\textwidth}{@{}P{0.5in}X@{}}
\toprule
\typed{R} & Retry: put things right first, such as closing the drive door
or turning the printer on, and try again.\\
\typed{A} & Abort: end the program that asked.\\
\typed{I} & Ignore: go on as if it had worked. Data may be lost; use it
only to rescue what can be read.\\
\bottomrule
\end{tabularx}

\begin{tabularx}{\textwidth}{@{}P{1.6in}X@{}}
\toprule
Not ready & The drive door is open or there is no diskette; the printer is
off or off-line.\\
Write protect & The diskette has a write-protect tab.\\
Seek, Sector not found, Data & The disk cannot be read there: the diskette
is damaged, or was written by a drive whose heads are out of line.\\
Non-DOS disk & The diskette is not formatted, or not for DOS.\\
No paper & The printer is out of paper.\\
\bottomrule
\end{tabularx}

\section{Common Messages}\ixs{messages}

\begin{longopts}
\lopt{Bad command or file name}{No internal command has that name, and no
  program or batch file of that name is in the current directory or the
  path.}
\lopt{File not found}{No file matches the name. Check the drive and path.}
\lopt{Invalid drive specification}{There is no such drive.}
\lopt{Invalid directory}{The directory named does not exist.}
\lopt{Invalid date}{Give the date as month, day and year, with the year in
  four digits: \typed{11-13-2026}.}
\lopt{Insufficient disk space}{The disk is full. Erase files you do not
  need, or use another disk.}
\lopt{File creation error}{The root directory is full, or a read-only file
  of that name is in the way.}
\lopt{Duplicate file name or File not found}{\typed{RENAME} found no such
  file, or a file with the new name exists.}
\lopt{Incorrect DOS version}{The command is from another version of DOS.
  Use the one from the RetrOS/LT diskette.}
\lopt{Out of environment space}{See ``The Environment'' in Section~5.}
\lopt{Insert diskette for drive B: and strike any key when ready}{The
  computer has one diskette drive, and a program asked for the other
  letter. Put in the diskette for that drive.}
\end{longopts}

\section{Programs That Check the DOS Version}\ixs{version, DOS}

RetrOS/LT tells programs it is DOS 2.10, so every program written for DOS
2.0 or 2.10 runs. A program written for DOS 3 or later that checks the
version says so and ends; one that does not check may use functions DOS 2
does not have, and fail. Use a version of the program made for DOS 2.

Some programs look for the system files by name, \file{IBMBIO.COM} and
\file{IBMDOS.COM}, to tell PC DOS from other systems; they do not find
them on a RetrOS/LT disk.

\section{Getting Files Back}

RetrOS/LT has no command to bring back an erased file; keep backup copies
with \typed{BACKUP}, \typed{COPY} or \typed{DISKCOPY}. If part of a file
cannot be read, \typed{RECOVER} rescues the rest. If \typed{CHKDSK} reports
lost clusters, \typed{CHKDSK /F} makes files of them, named
\file{FILE\var{nnnn}.CHK}, which you can look at with \typed{TYPE}.

\section{Reporting a Problem}

If you find that RetrOS/LT does something differently from PC DOS 2.10 on
the same computer, or wrongly, please write to The Lyon Mountain Company,
P.O. Box 290, Lyon Mountain, New York 12952 USA, or use the Reader's
Comment Form at the back of this book. Say what computer, memory and disks
you have, what you typed, and what appeared.
